\par\Needspace{6\baselineskip}\medskip
\hypertarget{native-085-article066}{}\pdfbookmark[2]{66.}{article-066}
\begin{center}66.\end{center}
E longitudine et latitudine puncti cuiusvis in sphaera coelesti eius rectascensio et declinatio derivantur per solutionem trianguli sphaerici, quod ab illo puncto polisque arcticis eclipticae et aequatoris formatur. Sit $\epsilon$ obliquitas eclipticae, $l$ longitudo, $b$ latitudo, $\alpha$ ascensio recta, $\delta$ declinatio, eruntque trianguli latera $\epsilon$, $90^0-b$, $90^0-\delta$; pro angulis lateri secundo et tertio oppositis accipere licebit $90^0+\alpha$, $90^0-l$ (siquidem trianguli sphaerici ideam maxima generalitate concipimus); angulum tertium lateri $\epsilon$ oppositum statuemus $=90^0-E$. Habebimus itaque per formulas art. 54
\SourceLeaf{086}
\[
\begin{aligned}
\sin(45^0-\tfrac12\delta)\sin\tfrac12(E+\alpha)&=\sin(45^0+\tfrac12 l)\sin(45^0-\tfrac12(\epsilon+b))\\[.5em]
\sin(45^0-\tfrac12\delta)\cos\tfrac12(E+\alpha)&=\cos(45^0+\tfrac12 l)\cos(45^0-\tfrac12(\epsilon-b))\\[.5em]
\cos(45^0-\tfrac12\delta)\sin\tfrac12(E-\alpha)&=\cos(45^0+\tfrac12 l)\sin(45^0-\tfrac12(\epsilon-b))\\[.5em]
\cos(45^0-\tfrac12\delta)\cos\tfrac12(E-\alpha)&=\sin(45^0+\tfrac12 l)\cos(45^0-\tfrac12(\epsilon+b)).
\end{aligned}
\]
\noindent Aequationes duae primae dabunt $\tfrac12(E+\alpha)$ atque $\sin(45^0-\tfrac12\delta)$; duae ultimae $\tfrac12(E-\alpha)$ atque $\cos(45^0-\tfrac12\delta)$; ex $\tfrac12(E+\alpha)$ et $\tfrac12(E-\alpha)$ habebitur $\alpha$ simulque $E$; ex $\sin(45^0-\tfrac12\delta)$ aut $\cos(45^0-\tfrac12\delta)$, quorum consensus calculo confirmando inserviet, determinabitur $45^0-\tfrac12\delta$ atque hinc $\delta$. Determinatio angulorum $\tfrac12(E+\alpha)$, $\tfrac12(E-\alpha)$ per tangentes suos ambiguitati non est obnoxia, quoniam tum sinus tum cosinus anguli $45^0-\tfrac12\delta$ positivus evadere debet.

Mutationes differentiales quantitatum $\alpha$, $\delta$ e mutationibus ipsarum $l$, $b$ secundum principia nota ita inveniuntur:
\[
\begin{aligned}
\mathrm{d}\alpha&=\frac{\sin E\cos b}{\cos\delta}\,\mathrm{d}l-\frac{\cos E}{\cos\delta}\,\mathrm{d}b\\[.6em]
\mathrm{d}\delta&=\cos E\cos b\,\mathrm{d}l+\sin E\,\mathrm{d}b.
\end{aligned}
\]
